% =========================================================================
% From End-State Equivalence to Per-Instruction Lockstep:
% A UVM Verification Environment for the Vortex RISC-V GPGPU
%
% IEEE conference format (IEEEtran). Compile on Overleaf or any TeXLive:
%   pdflatex vortex_uvm_paper.tex && pdflatex vortex_uvm_paper.tex
%
% All numbers and claims in this paper are sourced from the audited base
% docs/PAPER_BASE_EVALUATION.md (commit 3a7de25) -- every figure was
% re-verified against primary artifacts (UCDBs via vcover, RTL excerpts,
% committed investigation docs) on 2026-07-16.
% =========================================================================
\documentclass[conference]{IEEEtran}
\usepackage[utf8]{inputenc}
\usepackage[T1]{fontenc}
\usepackage{booktabs}
\usepackage{array}
\usepackage{url}
\usepackage{listings}
\usepackage{xcolor}
\usepackage{balance}

\lstset{
  basicstyle=\ttfamily\scriptsize,
  breaklines=true,
  frame=single,
  framerule=0.3pt,
  aboveskip=4pt, belowskip=4pt
}

\newcommand{\code}[1]{\texttt{\small #1}}

\begin{document}

\title{From End-State Equivalence to Per-Instruction Lockstep:\\
A UVM Verification Environment for the Vortex RISC-V GPGPU}

% ------------------------------------------------------------------
% Author block. Supervisor is currently in the Acknowledgment section;
% move into the author list if your venue/department requires it.
% ------------------------------------------------------------------
\author{
\IEEEauthorblockN{Samuel Moussa, Steven Ibrahim, Ahmad Sudky, Ahmad Fawzy, Abanoub Nabil}
\IEEEauthorblockA{Department of Electronics and Communications Engineering\\
Faculty of Engineering, Minia University, Egypt}
}

\maketitle

% =========================================================================
\begin{abstract}
General-purpose GPU (GPGPU) designs invert the assumptions of mainstream
processor verification: the device under test is a program-executing bus
\emph{master} whose stimulus is compiled kernels rather than bus
transactions, and whose SIMT execution model (warps, thread masks,
divergence, barriers) is outside the domain of standard RISC-V reference
models. We present a complete UVM verification environment for Vortex, an
open-source RISC-V GPGPU, built in two stages. The first stage is a
black-box environment whose single checker is byte-exact final-memory
equivalence against the project's functional simulator (SimX) integrated
over DPI-C; its distinguishing discipline is \emph{provable
non-vacuity}---permanent fault-injection tests that must stay red guard
every scoreboard change---together with machine-generated,
RTL-cited, per-configuration coverage exclusions and an explicitly
classified \emph{unverifiable} verdict class. This stage reaches 100\%
functional and 91.0\% total coverage on the primary configuration
(43/43 tests passing) and 85.2\% on a two-cluster configuration. The
second stage upgrades checking depth to an RVVI-style \emph{per-instruction
lockstep} against the stepping functional model---to our knowledge the
first for a SIMT GPGPU. We report five SIMT-specific alignment rules that
a scalar-CPU lockstep never encounters, a sound two-pass load-value feed
that renders racy, fenceless multi-core programs
instruction-granularity verifiable (residual~0 over 5{,}432 retirements on
the flagship case), a reference-model fetch bug found by the lockstep
itself, and a precisely characterized soundness boundary at asynchronous
interrupt timing.
\end{abstract}

\begin{IEEEkeywords}
GPGPU verification, UVM, RISC-V, SIMT, lockstep, RVVI, functional
coverage, fault injection, Vortex
\end{IEEEkeywords}

% =========================================================================
\section{Introduction}
\label{sec:intro}

Open-source RISC-V CPU cores enjoy a mature verification ecosystem:
constrained-random instruction generators \cite{riscvdv}, spec-complete
instruction-set simulators \cite{spike}, and step-and-compare frameworks
such as OpenHW \code{core-v-verif} with the RISC-V Verification Interface
(RVVI) \cite{corevverif,rvvi}. None of this transfers directly to an
open-source GPGPU. Vortex \cite{vortex-micro,vortex-carrv}, a RISC-V
GPGPU implementing RV32IMAF plus six SIMT control instructions
(\code{wspawn}, \code{tmc}, \code{split}, \code{join}, \code{bar},
\code{tex}), differs from a CPU DUT in three ways that shape the whole
verification problem:

\begin{enumerate}
\item \textbf{The DUT is a master.} Vortex initiates all memory traffic
      over AXI4; the testbench must \emph{respond}, not drive. Stimulus
      is therefore a compiled program image plus a launch sequence, not a
      stream of bus transactions.
\item \textbf{The reference-model domain is SIMT.} A scalar ISS such as
      Spike cannot execute warp-collective instructions, divergence
      stacks, or hardware barriers, so the natural golden model is the
      project's own cycle-level simulator, SimX---which was never
      designed to be a verification reference.
\item \textbf{Legal nondeterminism is first-class.} Fenceless multi-core
      programs have many legal outcomes; a golden model resolves shared
      memory ordering differently than timing-accurate RTL without either
      being wrong.
\end{enumerate}

This paper describes how we addressed all three, in two stages, and what
we believe generalizes. Our contributions are:

\begin{enumerate}
\item A complete black-box UVM environment for a SIMT GPGPU with a
      role-inverted agent architecture (DUT master, reactive slave
      agents backed by a memory model) and a DPI-integrated functional
      golden model (Section~\ref{sec:env}).
\item A \emph{trust-the-bench} methodology in which verdict non-vacuity
      is demonstrated, not assumed: permanent negative fault-injection
      tests (wrong-value and dropped-store) act as regression guards for
      the checker itself (Section~\ref{sec:gate0}).
\item An honest coverage-closure discipline: config-aware,
      machine-generated exclusions each citing the RTL line that proves
      structural unreachability; root-caused structural ceilings reported
      at true value rather than gamed; per-configuration coverage banks
      with a demonstration that cross-configuration merging is invalid
      (Section~\ref{sec:coverage}).
\item To our knowledge the first RVVI-style per-instruction lockstep for
      a SIMT GPGPU, including five SIMT-specific alignment rules derived
      from simulation evidence (Section~\ref{sec:lockstep}).
\item A sound two-pass load-value feed that makes racy fenceless
      multi-core programs verifiable at instruction granularity, with
      the residual mismatch count as the verdict
      (Section~\ref{sec:loadbus}); and a characterized soundness
      boundary at asynchronous interrupt delivery
      (Section~\ref{sec:boundary}).
\item Evidence that \emph{the golden model is itself a verification
      target}: the lockstep exposed and localized a reference-model
      instruction-fetch bug that had silently misclassified runs
      (Section~\ref{sec:refbug}).
\end{enumerate}

Throughout, we quantify an often-invisible cost: bring-up of the
DUT--reference co-simulation alone surfaced 43 distinct root-caused
issues before any meaningful verification could begin
(Section~\ref{sec:bringup}).

% =========================================================================
\section{Background and Related Work}
\label{sec:background}

\subsection{Vortex}
Vortex \cite{vortex-micro,vortex-carrv} is an open-source RISC-V GPGPU
organized as clusters of sockets of cores behind shared L2/L3 caches. Each
core is a five-stage in-order pipeline extended with SIMT machinery: a
warp scheduler, an IPDOM divergence stack, banked register files, and a
hardware barrier unit; execution units comprise ALU, FPU, LSU, SFU, and a
tensor-core unit (TCU). The memory system is a multi-banked, non-blocking,
\emph{write-through} cache hierarchy with per-bank MSHRs. Our work pins
one RTL commit and verifies two build configurations: a primary
single-cluster configuration (1 cluster / 1 core / 4 warps / 4 threads,
RV32, AXI) and a scale configuration (2 clusters / 2 cores / 4 warps / 4
threads).

\subsection{Reference models}
SimX is Vortex's C++ cycle-level simulator, intended for design-space
exploration. Crucially, its functional emulator steps
instruction-by-instruction, holds full SIMT architectural state, and is
cleanly decoupled from its timing model---properties we exploit in
Section~\ref{sec:lockstep}. SimX is written by the design team and is
therefore \emph{not independent} of the DUT; shared assumptions can hide
shared bugs. A scalar spec-authority ISS (Spike \cite{spike}) is planned
as a base-ISA cross-check but cannot serve as the primary golden model
because it cannot execute the SIMT instructions.

\subsection{Related verification frameworks}
OpenHW \code{core-v-verif} \cite{corevverif} established step-and-compare
verification of RISC-V CPU cores against a reference model over RVVI
\cite{rvvi}, commercialized in ImperasDV-style flows where
unpredictable inputs (loads from volatile memory, asynchronous
interrupts) are \emph{fed} from the DUT to the reference. riscv-dv
\cite{riscvdv} provides constrained-random RISC-V program generation. Our
environment began, deliberately, as an extension of the
\code{core-v-verif} CV32E40P environment toward a GPGPU target, and the
second stage of this work returns to that methodology family---but every
step-and-compare mechanism had to be rethought for SIMT semantics
(Section~\ref{sec:corrections}).

% =========================================================================
\section{Black-Box Environment Architecture}
\label{sec:env}

\subsection{Role-inverted agent architecture}
The environment comprises five UVM agents: host-control, DCR
(device-configuration register), and AXI agents are active; a custom
memory-protocol agent operates active-or-passive; a status agent is
passive. Because the DUT masters the bus, the AXI and memory agents are
\emph{reactive slaves}: drivers listen on AW/W/AR channels and serve B/R
responses by querying a sparse byte-addressed memory model that also holds
the pre-loaded program image. A virtual sequencer coordinates launch:
DCR writes (startup address, kernel arguments) while reset is held,
followed by liveness monitoring to a decoded-\code{EBREAK} completion
event.

Stimulus and tests are orthogonal by construction: a UVM \emph{test}
defines how to drive (configuration, launch protocol, fault-injection
modes); the \emph{program}---a compiled Vortex kernel ELF or a riscv-dv
assembly image---defines what executes. The two compose at the build
level, so every checker and coverage instrument applies uniformly to
directed and random stimulus.

\subsection{Golden-model integration}
SimX is compiled into the simulation as a DPI-C shared library. The
scoreboard drives it with the same program image and DCR configuration as
the RTL, runs it to completion, and performs a \emph{bidirectional}
end-state comparison: a forward pass checks every word the DUT wrote
against SimX, and a reverse pass checks words SimX wrote that the DUT
never wrote---catching dropped stores, which a forward-only comparison
is blind to. Sub-word stores are handled with per-byte validity masking;
read-only sections and language-runtime artifacts (\code{.got}) are
excluded with cited rationale. Both RTL and SimX are rebuilt per
configuration so reference and DUT always match topology, and
elaboration-time assertions compare every UVM topology parameter against
the RTL's macros, failing loudly on mismatch.

\subsection{Passive observability}
A passive probe is bound to each core's commit arbiter, exposing the full
retirement record (instruction UUID, warp ID, PC, thread mask,
destination register, per-lane data) for coverage and---in
stage two---for lockstep capture. Probes observe and never gate: the
black-box verdict comes only from state comparison against the golden
model.

\subsection{The cost of bring-up}
\label{sec:bringup}
Before the environment produced its first honest verdict, integration
surfaced \textbf{43 distinct root-caused issues} across three domains:
RTL simulation under the commercial simulator (macro type strictness,
protected diagnostics, third-party FPU packaging), the SimX DPI
integration (a RAM page-size constructor bug yielding zero-reads, C++
ABI/GLIBCXX mismatches against the simulator's bundled toolchain, global
state ordering hazards, exit-code protocol), and the UVM architecture
(dual-driver conflicts, dependency-ordered compilation). We report the
count because this integration cost is real, rarely published, and
dominated early schedule. Each issue is documented with symptom, root
cause, and fix in the project record.

% =========================================================================
\section{Trust the Bench: Provable Non-Vacuity}
\label{sec:gate0}

Three independent incidents during bring-up produced \emph{vacuous
passes}---green results that verified nothing:
(i)~the scoreboard mirrored every DUT write into SimX memory before
running it, so the reference computed from DUT-contaminated state and
comparisons passed \emph{regardless of RTL correctness};
(ii)~a smoke test passed having executed one instruction with zero
comparisons; and
(iii)~a hex-image base-address overflow left the code region empty, and
the DUT ``passed'' by fetching NOPs.

These incidents motivated a blocking methodology gate (``Gate-0''): no
coverage or regression number is reported until the bench itself is
proven. Its elements:

\begin{itemize}
\item \textbf{Permanent negative tests.} Two fault-injection tests ship
      with the environment: one flips a single bit in one DUT store; one
      drops a store entirely. The checker must catch both---the tests
      \emph{must stay red on injection}---and both run as regression
      guards after any scoreboard change. Injection later extends to the
      lockstep layer (Section~\ref{sec:lockstep}), where a single flipped
      bit must be caught at the exact UUID, PC, and SIMT lane.
\item \textbf{No fabricated metrics.} An instruction counter derived as
      \code{mem\_ops/3} was replaced with a true retired-instruction tap
      on the commit arbiter across all cores and lanes; an error gate
      that subtracted two from the UVM error count was removed;
      completion detection moved from an idle-threshold heuristic to
      decoded \code{EBREAK}.
\item \textbf{Derived, asserted parameters.} Interface widths are derived
      from the RTL package (a hand-coded ``ID width = 8'' turned out to
      be a 50-bit tag in reality) and checked by elaboration assertions.
\end{itemize}

Gate-0 also root-caused two startup-protocol findings in the RTL's
surrounding behavior: the core self-starts from reset while its DCR base
registers have no reset value---a race the environment closes by holding
reset until a configuration-done handshake---and a latent RTL bug in
which a watchdog macro computes \code{1**N}, which is identically~1, so a
stall timeout never scales with cache-hierarchy depth as intended.

% =========================================================================
\section{Stimulus and Honest Coverage Closure}
\label{sec:coverage}

\subsection{Stimulus}
The suite combines $\sim$30 directed kernels (real C kernels compiled
with the Vortex toolchain: vector arithmetic, all 13 FPU operation
classes, tensor-core WMMA, nested-divergence towers, barriers under
partial thread masks, MSHR saturation, a 232\,KB-text kernel, a 256\,KB
high-entropy store stress) with a working riscv-dv pipeline
(RV32IM target; machine-mode CSRs stripped and \code{ecall} rewritten for
the GPU target) contributing $\sim$12 random-program seeds, plus
plusarg-gated AXI slave stress modes (response throttling and read
flooding) that are proven byte-identical to default behavior when
disabled.

\subsection{Coverage model}
Seventeen covergroups span instruction classes per execution unit
(operation-decoded), SIMT divergence crossed with IPDOM depth,
warp/thread-mask crosses for collective operations, barrier and
warp-spawn behavior, a stall taxonomy crossed with windowed IPC, AXI
transaction fields, DCR/host activity, and system state. Each group
carries a written sufficiency rationale in a reviewed reference document.

\subsection{Closure discipline}
\label{sec:closure}
Three rules kept closure honest:

\begin{enumerate}
\item \textbf{Exclusions are structural, cited, and generated.} A
      generator script emits per-configuration exclusion files in which
      every waiver cites the RTL line proving unreachability---e.g.,
      write-response stability assertions are excluded because the AXI
      adapter hardwires \code{bready=1}; a barrier's global flag is
      excluded \emph{only} at single-core builds where it is
      structurally unreachable, and kept at $\ge$2 cores. The merge flow
      verifies zero ineffective exclusions.
\item \textbf{Ceilings are root-caused, not gamed.} Toggle coverage
      saturates near 78\%: the write-through cache never drives the
      512-bit line write-data fields, and PC/address high bits are
      constant for realistic program layouts. A maximum-entropy
      adversarial kernel moved aggregate toggle by $+0.02$\%,
      demonstrating the ceiling is structural. It is reported at true
      value; the founding toggle goal ($>$90\%) is reported as
      \emph{not met}.
\item \textbf{Per-configuration banks, never blended.} Merging coverage
      databases across configurations inflates the instance count
      (2{,}247\,$\rightarrow$\,8{,}260 between our two configurations)
      and mixes incompatible signal widths; per-configuration reporting
      is the only sound aggregation.
\end{enumerate}

\subsection{Results}
Table~\ref{tab:coverage} shows the banked results, reproduced from the
coverage databases with the vendor tool during the audit of this paper.
An observation of practical value: the tool's ``total'' is the
\emph{unweighted mean of seven category percentages} regardless of bin
count, so the lowest category is the largest lever---closing 11
directive bins moved the total more than 425{,}000 toggle bins could.

\begin{table}[t]
\caption{Coverage results per configuration (never blended). Reproduced
from the banked UCDBs via \code{vcover report -summary}.}
\label{tab:coverage}
\centering
\begin{tabular}{lrr}
\toprule
Metric & 1CL/1C/4W/4T & 2CL/2C/4W/4T \\
\midrule
Functional (covergroups) & \textbf{100.00\%} & 92.48\% \\
Statement   & 97.05\% & 96.19\% \\
Branch      & 91.16\% & 89.68\% \\
Condition   & 76.35\% & 69.57\% \\
Toggle      & 78.65\% & 74.25\% \\
Assertion   & 93.79\% & 73.96\% \\
Directive   & 100.00\% & 100.00\% \\
\midrule
\textbf{Total} & \textbf{91.00\%} & \textbf{85.16\%} \\
Tests passing  & 43/43 & 40/42 \\
\bottomrule
\end{tabular}
\end{table}

\subsection{The unverifiable verdict class}
\label{sec:unverifiable}
Runs the golden model cannot judge are \emph{classified}, never
force-compared: $\sim$10 random programs abort SimX on exotic encodings
(the current tree contains 56 \code{std::abort} sites in its
decoder/executor); and at two clusters, two fenceless/interrupt-heavy
random tests diverge from SimX deterministically \emph{per cluster}
(SimX's cluster-0 cores match the DUT byte-exactly while cluster-1
diverges by one propagating value)---root-caused by elimination to
unsynchronized cross-core memory ordering resolved differently by a
functional interleaving than by timing-accurate RTL, consistent with
Vortex's relaxed coherence \cite{vortex-micro}. Under end-state checking
these remained honest ``unverifiable'' entries. Retiring this bucket is
the strongest motivation for stage two.

% =========================================================================
\section{Per-Instruction Lockstep for SIMT}
\label{sec:lockstep}

\subsection{Architecture}
The functional emulator inside SimX already steps one instruction at a
time, returns an RVVI-shaped retirement record, and exposes per-lane
destination-register readback---it is most of an RVVI golden model,
needing only an export path. We serialize each retirement
\{UUID, core, warp, thread mask, PC, destination, per-lane value,
functional-unit type, volatile flag\} over the existing DPI bridge. On
the DUT side, the passive commit probe captures every retirement beat
under a runtime flag (default-off runs are proven byte-identical). A
dedicated lockstep scoreboard aligns the two streams and compares every
\emph{active SIMT lane}: PC, destination register, and written value,
with a four-way outcome taxonomy (matched, field mismatch, data
mismatch, orphan) and an end-of-test drain-empty check on both sides so
dropped or extra retirements cannot hide. Non-vacuity is again proven by
injection: a single flipped bit in one retirement is caught at the exact
UUID, PC, and lane.

On the primary configuration the proof-of-concept matched 1{,}035 of
1{,}035 architectural writebacks with zero orphans and zero mismatches.

\subsection{Five SIMT-specific alignment rules}
\label{sec:corrections}
Each of the following was forced by simulation evidence, and none arises
in a scalar-CPU lockstep. We consider them the transferable design rules
for SIMT step-and-compare:

\begin{enumerate}
\item \textbf{Retire order is not program order.} The commit arbiter
      retires whichever execution unit is ready (we observed an earlier
      UUID retiring after two later ones). The reference retires in
      program order. Alignment must key on the per-warp issue counter
      (UUID), sorted---never on position or cycle.
\item \textbf{The reference's UUID is not a cross-key.} SimX does not
      populate it; the practical key is per-(core, warp) program order,
      with the DUT's warp ID recovered from high UUID bits.
\item \textbf{One instruction is many retirement records.} SIMD-beat
      splits and, more subtly, \emph{partial-mask load writebacks}: the
      LSU commits lanes as memory responses arrive, so one \code{lw}
      appeared as two records with overlapping thread masks. Records
      must be aggregated by UUID with thread-mask union before
      comparison; any single in-flight slot keyed on start/end flags
      corrupts under interleaving.
\item \textbf{Load data is unobservable at the commit arbiter} (the LSU
      writeback path is asynchronous to it). Load values were first
      scoped to PC/ordering checks and later recovered soundly through a
      dedicated writeback probe and the feed of
      Section~\ref{sec:loadbus}.
\item \textbf{Performance-counter CSRs are model-divergent by
      definition} (cycle counts differ between timing RTL and a
      functional model). The reference flags them volatile and they are
      excluded from data comparison---exactly as commercial RVVI flows
      do.
\end{enumerate}

\subsection{Generalization across SIMT and topology}
With rules (1)--(5) in place and no further comparator changes, a
nested-divergence kernel matched 2{,}668/2{,}668 retirements through
split/join reconvergence. Multi-core attribution required no RTL change:
we proved from the RTL that the DUT's UUID embeds a flat global core
index matching the reference's core ID, and validated the comparator
lane-exactly across a configuration matrix spanning clusters
$\in\{1,2\}$, cores $\in\{1,2\}$, warps $\in\{2,4,8\}$, threads
$\in\{2,4\}$ (e.g., 3{,}333/3{,}333 retirements across four cores in two
clusters), all with zero mismatches and zero orphans.

\subsection{The reference model is a verification target}
\label{sec:refbug}
Running formerly-unverifiable programs under lockstep immediately paid
off \emph{against the golden model}: SimX fetched at the exact byte PC,
while the DUT's debug build preserves a full odd PC after an
indirect jump to an odd target but issues word-aligned fetches. SimX
therefore read misaligned bytes, decoded garbage, and aborted---runs
that were wrongly classified unverifiable. The fix word-aligns the
emulator fetch; notably, the ``obvious'' alternative of masking the JALR
target LSB was tried and \emph{reverted with evidence} (it desynchronizes
the reference from the DUT's odd PC, producing 21{,}968 phantom PC
mismatches). The decoder now prints the faulting PC and instruction word
before any abort, so every remaining unverifiable case self-documents.

% =========================================================================
\section{Verifying Racy Programs: a Two-Pass Load-Value Feed}
\label{sec:loadbus}

For fenceless multi-core programs, racy load values are legally
nondeterministic: the reference cannot predict them, and any interleaving
is architecturally valid. Commercial RVVI flows solve this by feeding
DUT-observed load data to the reference. Our implementation is a
\emph{sound two-pass trace replay}:

\begin{enumerate}
\item \textbf{Pass 1} runs DUT and reference independently and captures
      divergent loads keyed by (core, warp, PC, occurrence index)---
      UUIDs differ across the two models, so instruction identity is
      positional within this key.
\item \textbf{Feed.} Captured values are pushed to the reference over
      DPI; a \code{consumed==pushed} self-check guards feed alignment.
\item \textbf{Pass 2} re-runs the reference in-process with the feed
      active and re-compares everything. \emph{The residual mismatch
      count is the verdict}: pass-1 divergences are demoted to
      diagnostics only when the feed is armed, and any residual not
      explained by a fed racy load remains a hard error.
\end{enumerate}

On the pinned flagship case---a fenceless random program at two clusters
whose end-state divergence had been unexplainable for weeks under
end-state checking---lockstep first pinpointed the exact first divergent
instruction (a \code{mulhu} at PC \code{0x800004f4}, sequence 278,
cluster-1 cores only, with both register values in evidence; its inputs
had already diverged, proving an upstream shared load rather than a
compute bug). The feed then captured 20 racy loads, whose replay
collapsed a 138-divergence cascade to \textbf{residual 0 over
5{,}432/5{,}432 retirements}, and the deferred end-state comparison
passed. A run that was ``unverifiable'' is now \emph{positively verified
at instruction granularity}, with no waiver. Default-path regressions
remain byte-identical and both stage-one injection guards remain red.

% =========================================================================
\section{Soundness Boundary: Interrupt Timing}
\label{sec:boundary}

The same feed applied to an interrupt-heavy random test at two clusters
collapses 116 divergences to 7---and provably cannot reach zero. We
re-keyed the feed from a global ordinal to (core, warp, PC, occurrence)
and obtained a \emph{bit-identical residual of 7}, disproving the
hypothesis that the residual is a feed-alignment artifact. The residual
is genuine interrupt-\emph{delivery} timing: the timing-accurate DUT and
the functional reference take the interrupt at different instruction
boundaries, so the interrupt-affected path executes different iteration
counts, and no amount of load-\emph{data} feeding can align \emph{when}
an asynchronous event fires. The end-state comparison passes, so the run
is dispositioned end-state-verified with a classified
instruction-granularity residual---not forced green.

This yields a precise statement of the method's domain: \textbf{two-pass
trace replay is a fixed point for data-only divergence; asynchronous-input
timing requires a step-follower reference}, i.e., stepping the reference
one instruction at a time as a follower of the DUT retirement stream with
interrupt delivery injected at the same retired-instruction boundary---the
ImperasDV-style architecture, which we have scoped as future work.

% =========================================================================
\section{Limitations}
\label{sec:limitations}

(i)~SimX is written by the design team and is not independent of the DUT;
a Spike base-ISA audit is planned but not yet layered in.
(ii)~$\sim$10 random-program runs remain unverifiable pending
abort-hardening of the reference decoder (56 abort sites measured at the
time of writing); each fix shrinks the bucket.
(iii)~The founding toggle-coverage goal is unmet for root-caused
structural reasons (Section~\ref{sec:closure}).
(iv)~Interrupt-timing tests are end-state-verified only
(Section~\ref{sec:boundary}).
(v)~The raw simulation logs behind the lockstep configuration-matrix
tallies were session artifacts; the tallies are recorded in committed
project documents, while coverage, investigation, and commit evidence
remain reproducible from in-tree artifacts.

% =========================================================================
\section{Future Work}
\label{sec:future}

A single-pass \emph{step-follower} lockstep with interrupt-delivery
alignment (closing Section~\ref{sec:boundary}'s residual); the Spike
independence audit; harvesting the RTL's native runtime assertions into
the UVM error gate; formal proof of the structural-unreachability
exclusions on the cache arbiter and MSHR allocator (upgrading cited
waivers to proofs); a UVM register-abstraction layer for the DCR block;
and CI-driven parallel regression with automated, traceable sign-off.

% =========================================================================
\section{Conclusion}
\label{sec:conclusion}

We presented a two-stage verification program for an open-source SIMT
GPGPU. The black-box stage demonstrates that a program-executing
master DUT can be verified rigorously with a functional-simulator golden
model---\emph{provided} the bench itself is proven non-vacuous by
standing fault injection, coverage closure is disciplined by RTL-cited
structural evidence, and everything the reference cannot judge is
classified rather than forced. The lockstep stage shows that RVVI-class
per-instruction checking extends to SIMT execution once five specific
alignment rules are observed; that a sound two-pass load feed makes racy
multi-core programs verifiable at instruction granularity; and that the
approach's boundary at asynchronous interrupt timing can be stated---and
proven---precisely. Along the way the lockstep found and fixed a bug in
the reference model itself, reinforcing the campaign's central lesson:
\emph{verify the verifier}.

% =========================================================================
\section*{Acknowledgment}
This work was carried out as a graduation project at the Faculty of
Engineering, Minia University, under the supervision of
Dr.~Hossam~Hassan, whose guidance is gratefully acknowledged.

% =========================================================================
\balance
\begin{thebibliography}{9}

\bibitem{vortex-micro}
B.~Tine, K.~P. Yalamarthy, F.~Elsabbagh, and H.~Kim,
``Vortex: Extending the RISC-V ISA for GPGPU and 3D-graphics,''
in \emph{Proc. 54th IEEE/ACM Int. Symp. Microarchitecture (MICRO-54)},
2021, pp. 754--766.

\bibitem{vortex-carrv}
F.~Elsabbagh, B.~Tine, P.~Roshan, E.~Lyons, E.~Kim, D.~E. Shim,
L.~Zhu, S.~K. Lim, and H.~Kim,
``Vortex: OpenCL compatible RISC-V GPGPU,''
in \emph{Workshop on Computer Architecture Research with RISC-V (CARRV)},
2019.

\bibitem{corevverif}
OpenHW Group, ``core-v-verif: functional verification project for the
CORE-V family of RISC-V cores,'' [Online]. Available:
\url{https://github.com/openhwgroup/core-v-verif}

\bibitem{rvvi}
OpenHW Group / Imperas, ``RVVI: RISC-V Verification Interface,''
[Online]. Available: \url{https://github.com/riscv-verification/RVVI}

\bibitem{riscvdv}
CHIPS Alliance, ``riscv-dv: SV/UVM based instruction generator for
RISC-V processor verification,'' [Online]. Available:
\url{https://github.com/chipsalliance/riscv-dv}

\bibitem{spike}
RISC-V International, ``Spike, the RISC-V ISA simulator,'' [Online].
Available: \url{https://github.com/riscv-software-src/riscv-isa-sim}

\bibitem{uvm}
\emph{IEEE Standard for Universal Verification Methodology Language
Reference Manual}, IEEE Std 1800.2-2020, 2020.

\end{thebibliography}

\end{document}
